\chapter{Complexity — Codes2 Tape}
%%% generated by the blueprint pipeline from TCSlib.Complexity.TuringMachine.Codes2Tape
\section{Overview}

Declarations of the Lean module \texttt{TCSlib.Complexity.TuringMachine.Codes2Tape}.

\section{Declarations}

\begin{definition}[Finite-state two-tape machine code]
\lean{Turing.Code2TM}
\label{Turing.Code2TM}
\leanok
\uses{Turing.MultiTapeTM}
A code for a deterministic machine with two work tapes and Boolean tape symbols consists
of a natural number $n$ and such a machine whose state set is the finite set
$0,1,\ldots,n$ (that is, a set of $n+1$ states).
\end{definition}

\begin{definition}[Finite-machine representation]
\lean{Turing.Code2TM.toFinTM}
\label{Turing.Code2TM.toFinTM}
\leanok
\uses{Turing.Code2TM, Turing.FinTM}
Given a finite-state two-tape machine code with $n+1$ states, regard its underlying
machine as a finite Turing machine over the Boolean alphabet, with exactly two work
tapes and state set $0,1,\ldots,n$.
\end{definition}

\begin{definition}[Binary serialization of a machine]
\lean{Turing.Code2TM.serialize}
\label{Turing.Code2TM.serialize}
\leanok
\uses{Turing.Code2TM, Turing.actionBits₂, Turing.pairEncode, Turing.unaryFin, Turing.workPair}
For a finite-state two-tape machine code with $n+1$ states, define its binary
serialization by encoding $n$ in binary and pairing that code with a binary list
consisting of the initial state followed by the encoded transition actions. The
transition table is listed in increasing state order, and for each state in the order of
the three possible input symbols (blank, false, true), the three possible symbols on the
first work tape, and the three possible symbols on the second work tape.
\end{definition}

\begin{definition}[Machine coding scheme]
\lean{Turing.MachineCode2}
\label{Turing.MachineCode2}
\leanok
\uses{Turing.Code2TM}
A machine coding scheme consists of an encoder from finite-state two-tape machine codes
to finite Boolean strings, a decoder from every finite Boolean string to a finite-state
two-tape machine code, and the property that decoding the encoding of any machine
followed by any finite number of true bits recovers that machine.
\end{definition}

\begin{theorem}[Decoding an unpadded two-tape machine code]
\lean{Turing.MachineCode2.decode_encode}
\label{Turing.MachineCode2.decode_encode}
\leanok
\uses{Turing.Code2TM, Turing.MachineCode2}
\difficulty{1}
Let a representation scheme encode machines with two work tapes, binary alphabet, and a
finite nonempty state space of the specified form as bit strings, and decode every bit
string to such a machine. Suppose that, for every such machine and every natural number
of padding bits, decoding its code followed by that many $1$ bits returns the machine.
Then decoding the unpadded code of each such machine returns that machine.
\end{theorem}

\begin{definition}[Effective machine coding scheme]
\lean{Turing.EffectiveMachineCode2}
\label{Turing.EffectiveMachineCode2}
\leanok
\uses{Turing.Code2TM.serialize, Turing.FinTM, Turing.FinTM.ComputesFunInTime, Turing.MachineCode2}
An effective machine coding scheme consists of a machine coding scheme, a finite Turing
machine over the Boolean alphabet, a natural-valued time-bound function on
natural-number inputs, and the property that this machine computes the serialization of
the decoded input within that time bound on every input.
\end{definition}

\begin{definition}[Halting fallback machine]
\lean{Turing.zfBFallback}
\label{Turing.zfBFallback}
\leanok
\uses{Turing.Code2TM}
The fallback code represents a two-work-tape deterministic Turing machine with one
state, whose transition always halts immediately: it leaves the input and work-tape
heads stationary, writes nothing, and emits no output.
\end{definition}

\begin{lemma}[Fallback serialization length]
\lean{Turing.zfBFallback_serialize_length}
\label{Turing.zfBFallback_serialize_length}
\leanok
\uses{Turing.Code2TM.serialize, Turing.actionBits₂, Turing.length_pairEncode, Turing.optBoolBits, Turing.optOptBoolBits, Turing.optStateBits, Turing.signBits, Turing.unaryFin, Turing.zfBFallback}
\difficulty{1}
The serialization of the fixed fallback machine has length $354$.
\end{lemma}

\begin{definition}[Two-tape machine code format]
\lean{Turing.twoTapeFormat}
\label{Turing.twoTapeFormat}
\leanok
\uses{Turing.Code2TM, Turing.Code2TM.serialize, Turing.CodeFormat, Turing.CodeTable2, Turing.actionBits₂_eq_actionBitsK, Turing.codeEraseTable2, Turing.codeReadTable2, Turing.codeReadTable2_append, Turing.codeReadTable2_sound, Turing.codeSkipTable2, Turing.codeTable2Bits, Turing.codeTable2Bits_length, Turing.workPair, Turing.zfBFallback}
There is a code format for deterministic machines with two work tapes and binary tape
symbols. A code records its number of states, a designated initial state, and a
transition table; its serialization pairs the binary expansion of the state-count
parameter with the unary initial-state index followed by the serialized table. The
format has a table reader and an erased table skipper, a guard constant of $351$ bits
per state, and the halting fallback machine described above; assembling a state count,
initial state, and table recovers every code, and the reader and skipper satisfy the
stated consistency and length properties.
\end{definition}

\begin{definition}[Total decoding of two-tape codes]
\lean{Turing.zfBDecode}
\label{Turing.zfBDecode}
\leanok
\uses{Turing.Code2TM, Turing.CodeFormat.decode, Turing.twoTapeFormat}
Every finite binary string is decoded as a code for a deterministic two-work-tape Turing
machine; strings that do not form a valid code are interpreted as the format's fallback
machine.
\end{definition}

\begin{theorem}[Existence of a two-tape effective code scheme]
\lean{Turing.exists_effectiveMachineCode2}
\label{Turing.exists_effectiveMachineCode2}
\leanok
\uses{Turing.Code2TM.serialize, Turing.CodeFormat.decode_serialize_pad, Turing.CodeFormat.exists_canonizer, Turing.EffectiveMachineCode2, Turing.codePrimSkipTable2, Turing.twoTapeFormat, Turing.zfBDecode}
\difficulty{4}
There exists an effective coding scheme for deterministic two-work-tape Turing machines
over the binary alphabet: it has an encoding and a total decoder such that, for every
coded machine $M$ and every natural number $m$, decoding the encoding of $M$ followed by
$m$ copies of the bit $1$ returns $M$. The scheme also includes a finite binary-alphabet
Turing machine and a time bound that, on every binary input string, compute the fixed
serialization of the machine decoded from that string.
\end{theorem}

\begin{definition}[Uniform machine-code system]
\lean{Turing.UniformMachineCode2}
\label{Turing.UniformMachineCode2}
\leanok
\uses{Turing.Code2TM.toFinTM, Turing.EffectiveMachineCode2, Turing.FinTM, Turing.FinTM.ComputesInTime, Turing.pairEncode}
A uniform machine-code system consists of an effective code for deterministic
two-work-tape machines over binary symbols, together with a finite binary-input
simulator and a natural-number degree. For every code, input, and time bound, the
simulator decides whether the decoded machine outputs the one-bit word $[1]$ on that
input within the given bound: in the affirmative case it outputs $[1]$, and otherwise it
outputs $[0]$. In either case its computation takes exactly the prescribed time bound
$d(|\alpha|+|x|+t+1)^d$, where $d$ is the degree and $\alpha,x,t$ are the code, input,
and bound; the code system also includes a finite machine computing the serialization of
each decoded machine, with some time bound depending on input length.
\end{definition}

\begin{lemma}[Agreement of machine encodings]
\lean{Turing.Code2TM.serialize_eq_serializeK}
\label{Turing.Code2TM.serialize_eq_serializeK}
\leanok
\uses{Turing.Code2TM, Turing.Code2TM.serialize, Turing.actionBits₂_eq_actionBitsK, Turing.codeTable2Bits, Turing.serializeK, Turing.tableBitsK_two}
\difficulty{3}
For every two-work-tape deterministic Turing machine $M$ with Boolean alphabet and a
finite state set, the serialization of $M$ as a two-work-tape machine equals its
serialization using the general $k$-tape encoding with $k=2$.
\end{lemma}

\begin{definition}[Two-tape machine format]
\lean{Turing.twoTapeMachineFormat}
\label{Turing.twoTapeMachineFormat}
\leanok
\uses{Turing.Code2TM.serialize_eq_serializeK, Turing.MachineFormat, Turing.codeSkipTable2_eq_codeSkipTableK, Turing.twoTapeFormat}
There is a machine format whose codes represent deterministic machines with exactly two
work tapes, binary tape symbols, and state set $\{0,\ldots,n\}$ when the code's
state-count parameter is $n$. Each code's machine serialization is the fixed two-tape
serialization, and the format's table skipper is the fixed two-tape table skipper.
\end{definition}

\begin{theorem}[Existence of a two-tape uniform machine code]
\lean{Turing.exists_uniformMachineCode2}
\label{Turing.exists_uniformMachineCode2}
\leanok
\uses{Turing.Code2TM.serialize, Turing.CodeFormat.decode, Turing.CodeFormat.decode_serialize_pad, Turing.MachineFormat.exists_polyCanonizer, Turing.MachineFormat.exists_uniformSimulator, Turing.UniformMachineCode2, Turing.twoTapeFormat, Turing.twoTapeMachineFormat}
\difficulty{4}
There exists a code for deterministic Boolean machines with two work tapes, consisting
of an encoder and a total decoder such that decoding an encoded machine followed by any
number of $1$ bits returns that machine. The code also includes a finite-state Boolean
machine that computes the canonical serialization of each decoded machine within its
specified time bound, and a finite-state Boolean simulator with a natural number $d$.
For every program string $\alpha$, input string $x$, and time bound $t$, if the decoded
machine computes output $[1]$ on $x$ at exactly time $t$, the simulator computes $[1]$
on $\langle\langle\mathrm{bits}(t),\alpha\rangle,x\rangle$ at exactly
$d(|\alpha|+|x|+t+1)^d$ steps; if it does not compute $[1]$ at exactly time $t$, the
simulator computes $[0]$ on that input at exactly the same number of steps.
\end{theorem}
